\section{Proposed Framework Architecture}
\label{sec:framework}

The overall architectural design of \textbf{MultiAgentTopologyEval} is depicted in Figure~\ref{fig:system_architecture}. The framework is organized into six modular subsystems: (1) Benchmark Task Loader and Anti-Leakage Sanitizer, (2) Topology Orchestration Hub, (3) Role-Specialized Agent Core, (4) Universal LLM Provider Layer, (5) API Safety and Guardrails Engine, and (6) Hybrid Evaluation and Statistical Analytics Pipeline.

\begin{figure*}[t]
\centering
\begin{tikzpicture}[
    auto,
    block/.style={rectangle, draw=blue!80!black, fill=blue!8, thick, rounded corners=4pt, minimum width=2.4cm, minimum height=1.6cm, align=center, font=\scriptsize},
    subblock/.style={rectangle, draw=gray!70, fill=white, rounded corners=2pt, minimum width=2.1cm, minimum height=0.4cm, align=center, font=\tiny},
    box/.style={rectangle, draw=black!70, thick, rounded corners=5pt, inner sep=6pt, fill=gray!4},
    boxtitle/.style={font=\scriptsize\bfseries, text=black!90},
    arrow/.style={-{Stealth[length=2.2mm]}, thick, draw=blue!75!black},
    darrow/.style={{Stealth[length=2.0mm]}-{Stealth[length=2.0mm]}, thick, draw=teal!80!black}
]

    % Stage 1: Benchmark Suite
    \node[block, fill=orange!10, draw=orange!80!black] (stage1) {
        \textbf{Benchmark Suite}\\[2pt]
        \tiny 35 Canonical Tasks\\
        \tiny 5 Task Categories\\
        \tiny Explicit Constraints $\mathcal{C}$\\
        \vspace{2pt}
        \fcolorbox{orange!80!black}{white}{\tiny Anti-Leakage Isolation}
    };

    % Stage 2: Topology Orchestrator
    \node[block, right=0.9cm of stage1, fill=blue!10, draw=blue!80!black] (stage2) {
        \textbf{Topology Hub}\\[2pt]
        \tiny \textsc{Single} (1-turn)\\
        \tiny \textsc{Chain} (Linear SOP)\\
        \tiny \textsc{Star} (Hub-Spoke)\\
        \tiny \textsc{Tree} (Hierarchical)\\
        \tiny \textsc{Graph} (Mesh)
    };

    % Stage 3: Agent Core
    \node[block, right=0.9cm of stage2, fill=teal!10, draw=teal!80!black] (stage3) {
        \textbf{Agent Core}\\[2pt]
        \tiny BaseAgent Class\\
        \tiny Planner $\cdot$ Researcher\\
        \tiny Analyst $\cdot$ Critic\\
        \tiny Finalizer\\
        \tiny Memory \& State Log
    };

    % Stage 4: Provider Abstraction
    \node[block, right=0.9cm of stage3, fill=purple!10, draw=purple!80!black] (stage4) {
        \textbf{LLM Provider Layer}\\[2pt]
        \tiny Universal Adapter\\
        \tiny Google Gemini 3.1\\
        \tiny OpenAI GPT-4o\\
        \tiny Anthropic Claude 3.5\\
        \tiny Deterministic Mock
    };

    % Stage 5: Guardrails
    \node[block, below=0.8cm of stage3, fill=red!8, draw=red!80!black] (stage5) {
        \textbf{Guardrails \& Safety}\\[2pt]
        \tiny Cost Caps (\$0.50/run)\\
        \tiny Step Limit ($K \le 20$)\\
        \tiny Retry \& Backoff\\
        \tiny Error Taxonomy Log
    };

    % Stage 6: Evaluation Pipeline
    \node[block, right=0.9cm of stage4, fill=green!10, draw=green!70!black] (stage6) {
        \textbf{Hybrid Evaluator}\\[2pt]
        \tiny Constraint Engine\\
        \tiny Exact / Regex Matching\\
        \tiny Token \& Cost Profiling\\
        \tiny Latency \& Failure Audits\\
        \tiny Statistical Tests
    };

    % Flow Connections
    \draw[arrow] (stage1) -- node[above, font=\tiny] {Task $\langle p, \mathcal{C} \rangle$} (stage2);
    \draw[arrow] (stage2) -- node[above, font=\tiny] {Instantiate} (stage3);
    \draw[darrow] (stage3) -- node[above, font=\tiny] {Inference} (stage4);
    \draw[darrow] (stage3) -- node[right, font=\tiny] {Tracking} (stage5);
    \draw[darrow] (stage4) |- (stage5);
    \draw[arrow] (stage4) -- node[above, font=\tiny] {Traces \& $\hat{y}$} (stage6);
    
    % Evaluation Feedback / Reference Answer
    \draw[arrow, dashed, draw=orange!80!black] (stage1.south) -- ++(0,-1.4) -| node[near start, below, font=\tiny] {Ground-Truth Quarantined References ($y^*, \mathcal{P}$)} (stage6.south);

\end{tikzpicture}


\caption{System Architecture of MultiAgentTopologyEval: Decomposing execution across benchmark loading, anti-leakage isolation, topology orchestration, universal model provider abstraction, cost guardrails, and hybrid deterministic evaluation.}
\label{fig:system_architecture}
\end{figure*}

\subsection{Benchmark Loader and Anti-Leakage Protection}
The Benchmark Loader (\texttt{app/benchmark/loader.py}) ingests structured task suites from \texttt{data/tasks/benchmark\_tasks.json}. To maintain rigorous academic validity and eliminate data contamination, the framework enforces strict \textit{anti-leakage isolation} between agent execution and evaluation. When a task is dispatched to an architecture via \texttt{TaskDefinition.to\_agent\_prompt\_dict()}, ground-truth fields (\texttt{reference\_answer} and \texttt{expected\_properties}) are stripped from the payload. Agents receive only the operational prompt $p$, the task category, and explicit constraints $\mathcal{C}$. Authoritative reference answers remain strictly quarantined within the downstream evaluation engine.

\subsection{Role-Specialized Agent Core}
Agents inherit from a unified base class (\texttt{BaseAgent}) that encapsulates state representation, message queues, execution telemetry, and provider invocation. The framework instantiates five specialized roles:
\begin{enumerate}
    \item \textbf{Planner Agent:} Parses complex objectives, decomposes workflows into structured sub-tasks, and formulates step-by-step milestones.
    \item \textbf{Researcher Agent:} Performs domain knowledge exploration, information retrieval, and architectural context synthesis.
    \item \textbf{Analyst Agent:} Conducts quantitative constraint verification, algorithmic trade-off evaluation, and rule compliance checking.
    \item \textbf{Critic Agent:} Reviews intermediate proposals, identifies edge-case vulnerabilities, and provides constructive refinement critiques.
    \item \textbf{Finalizer Agent:} Consolidates multi-agent deliberations, reconciles conflicting peer outputs, and formats the authoritative final response.
\end{enumerate}

Every inter-agent communication event is structured as an immutable \texttt{AgentMessage} object logging the \texttt{sender\_id}, \texttt{receiver\_id}, \texttt{timestamp}, \texttt{message\_type}, and raw textual content.

\subsection{Universal Provider Abstraction Layer}
To isolate topology behavior from specific vendor APIs, MultiAgentTopologyEval implements an abstract provider interface (\texttt{LLMProvider}) that standardizes inference execution across heterogeneous cloud and local models:
\begin{itemize}
    \item \textbf{Standardized Output Schema:} Every API call returns an \texttt{LLMResponse} dataclass encapsulating the generated content, model identifier, wall-clock latency ($\pm 0.1\,$ms), decoupled input and output tokens, and calculated USD expenditure.
    \item \textbf{Multi-Backend Support:} Native integrations are provided for Groq (\texttt{llama-3.3-70b-versatile}), Google Gemini (\texttt{gemini-3.1-flash-lite}), OpenAI (\texttt{gpt-4o-mini}), Cohere (\texttt{command-a-plus}), Cerebras, and DeepSeek.
    \item \textbf{Deterministic Offline Mock Mode:} A built-in \texttt{MockProvider} generates deterministic, rule-based response tokens with zero external API calls. This enables 100\% offline unit testing, CI/CD regression verification, and student demonstrations at zero financial cost.
    \item \textbf{Safe Connectivity Testing:} Provider instances implement a non-destructive ping protocol (\texttt{test\_connection()}) that verifies network endpoints and authentication status while rigorously redacting credentials from log files.
\end{itemize}

\subsection{Resource Guardrails and Failure Trapping}
Unconstrained multi-agent loops can rapidly exhaust API budgets or enter infinite circular conversations. MultiAgentTopologyEval incorporates a proactive guardrail monitor (\texttt{GuardrailTracker}):
\begin{itemize}
    \item \textit{Call Caps:} Enforces strict thresholds on total API calls per task ($N_{\text{task}} \le 15$) and per experiment batch ($N_{\text{exp}} \le 60$).
    \item \textit{Loop Detection:} Detects circular message exchange by trapping agent turns exceeding maximum allowed iterations ($K_{\text{iter}} \le 5$).
    \item \textit{Pre-Flight Dry Run:} Pre-calculates expected token consumption and financial expenditure prior to live API dispatch.
    \item \textit{Fault Handling:} Intercepts and categorizes provider-level exceptions into structured failure modes (HTTP 429 rate limits, request timeouts, malformed JSON schemas, and constraint breaches).
\end{itemize}

